% ============================================================
%  paper-export :: BRIEF 预设的落款（经 --include-after-body 注入）
%
%  通报/纪要的署名在正文末尾右下角，不在封面：
%
%              微航科技有限公司
%              编制：禹贵祥
%              2026 年 8 月 5 日
%
%  距右边界留 4 个字宽 —— 公文里落款不顶到版心右边缘。
%  三行都可缺省；全缺省时整块不输出，不留空白。
% ============================================================

\makeatletter

\providecommand{\peorg}{}
\providecommand{\peauthorline}{}
\providecommand{\pesignoffdate}{}

\newcommand{\pe@sig@ifempty}[3]{\ifx#1\@empty#2\else#3\fi}

% 三个字段全空就整块跳过
\newif\ifpe@hassig \pe@hassigfalse
\ifx\peorg\@empty\else\pe@hassigtrue\fi
\ifx\peauthorline\@empty\else\pe@hassigtrue\fi
\ifx\pesignoffdate\@empty\else\pe@hassigtrue\fi

\ifpe@hassig
  % 可压缩的间距：正文离页底还宽裕时留足 1.6 行，页面紧张时一路压到 0.2 行，
  % 尽量让落款跟正文同页 —— 固定值会为了三行落款白开一页。
  \par\vspace{1.6\baselineskip plus 0pt minus 1.4\baselineskip}
  % flushright 会重置 \parindent，落款不需要缩进；右侧 4em 由 \hspace 补
  \begingroup
    \fangsong\zihao{3}\linespread{1.35}\selectfont
    \begin{flushright}
      % \hspace* 而非 \hspace：行尾的普通 glue 会在断行处被丢弃，
      % 带星号的版本不可丢弃，右侧 4 字宽才真的留得住。
      \pe@sig@ifempty{\peorg}{}{\peorg\hspace*{4em}\par}
      \pe@sig@ifempty{\peauthorline}{}{编制：\peauthorline\hspace*{4em}\par}
      \pe@sig@ifempty{\pesignoffdate}{}{\pesignoffdate\hspace*{4em}\par}
    \end{flushright}
  \endgroup
\fi

\makeatother
